\section{The actual coarse lower rank and density}
\label{sec:rank-lower}

\status{The weighted lower count, the actual kept-block injection and the
three source loss factors are proved under literal $\mathrm{Std}(H=4)$
and $C_0\ge880$. The shared identities RU1--RU3 identify this count with
the actual lattice rank. With RU8 this gives the source two-sided
half-density window. Uniform parameter supply is a separate input.}

Write $r=X.m$, $k=r+1$, $\lambda=C_0N_1^2\log a$,
$\epsilon=N_1\lambda/k$, and $\delta=X.\mathrm{de}$.
Fix an actual segment and compatible digits, and put
\[
 P_i=p^{G-d_i},\qquad A_*=\left\lfloor\frac{a-E_0}{p^G}\right\rfloor,
 \qquad
 A_i=\frac{b_i-\sum_c r^z_{ic}}{P_i},\quad
 C_i=\frac{a_i^{\rm top}-\sum_c r^y_{ic}}{P_i},\quad B=b_+.
\]
The quotients are nonnegative integers by compatibility.
Here $0<d_0<\cdots<d_{s-1}<G$ are odd and $G$ is even.
Let
$S_d(g)=\binom{g+d-1}{d-1}$ and
\[
 \mathcal D(A)=S_k(A)\prod_{i<s}S_k(A_i)S_k(C_i),\qquad
 \mathcal B=\mathcal D(A_*)S_k(B)=\mathrm{mmB}.
\]
For a real margin $h\ge0$, put $M=\lceil h\rceil$ and suppose
$M\le\epsilon A_*$. The claimed exact lower bound is
\[
 \mathrm{rankOnDigits}(h)\ \ge\
 (1-e^{-\lambda/8})(1-e^{-\lambda/3})
 \left(1-\frac{rM}{A_*}\right)\mathcal B.
 \tag{RD.1}
\]

\paragraph{The actual pair distribution.}
Put $A=A_*-M$. A pair consists of a signed Pieri chain
$\mu^0,\ldots,\mu^{2s+1}$ in $k$ coordinates, starting at zero,
with initial vector degree $A$ and subsequent pairs of signed degrees
$(+A_i,-C_i)_{0\le i<s}$, and an $r$-coordinate weight $\ell$
interlacing the endpoint.
The weight of a pair is the Weyl dimension $w_r(\ell)$.
The signed path identity and the interlacing branching identity,
proved in RU2--RU3, give
\[
                       \sum_{\text{pairs}}w_r(\ell)=\mathcal D(A).
 \tag{RD.2}
\]
Delete the first coordinate, writing $K_t^n=\mu_{t+1}^n$, and set
\[
\begin{gathered}
 c_i=\mu_0^{2i+1},\quad \rho_i=K_0^{2i+2},\quad
 W_i=\sum K^{2i+2}-\sum K^{2i+1},\\
 O_i=c_i+A_i-c_{i+1},\quad S_i=O_i-W_i,\quad
 u_i=O_i-\rho_i,\quad n_i=-\ell_{r-s+i}.
\end{gathered}
\]
Horizontal interlacing gives $0\le W_i\le A_i$, $0\le S_i\le C_i$
and $\rho_i\le c_i,c_{i+1}$. A pair is \emph{kept} precisely when
\[
 u_i+\frac{4(A_i+r)}{p^2}
       \le\frac\lambda k\max(c_i,A_i),\qquad
 C_i(1-\lambda/k)\le n_i\quad(i<s).
 \tag{RD.3}
\]

\paragraph{Why the bad weighted mass is small.}
The standard hypotheses and odd spacing imply
\[
\begin{gathered}
 k\ge10,\quad 4\lambda\le k,\quad p\ge k^2,\quad
 A,A_i,C_i\ge k^3,\quad 2s+1\le r,\\
 \sum_{l<i}A_l\le\frac{4A_i}{p^2},\qquad
 \sum_{l<i}C_l\le\frac{4C_i}{p^2}.
\end{gathered}
\]
For the geometric bounds, the consecutive degrees satisfy
$p^2A_i\le2A_{i+1}$ and $p^2C_i\le2C_{i+1}$: multiply by
$P_{i+1}$, use $P_i\ge p^2P_{i+1}$, and use the degree windows and
$kP_{i+1}\le b'/25$.
Following the upper and lower interlacing boundaries yields
\[
\begin{gathered}
 0\le\rho_i\le A_i(1+4/p^2),\quad u_i\ge-4A_i/p^2,
 \quad0\le n_i\le C_i(1+4/p^2),\\
 0\le\ell_0\le c_s,\quad
 0\le\ell_t\le\rho_{s-t}\ (1\le t\le s),\quad
 \ell_t=0\ (s<t<r-s).
\end{gathered}
\]
For example, negative mass of a tail grows by at most $C_i$ at a
covector step and decreases at a vector step; following its last
coordinate bounds $n_i$ by $\sum_{l\le i}C_l$.

Here is the quantitative core. Put $q=4/p^2$, $\eta=k^{-2}$,
$\tau=\lambda/(4k)$ and
\[
 \widehat c_i=c_i+r,\quad \widehat A_i=(A_i+r)(1+q),\quad
 \widehat\rho_i=\rho_i+r,\quad\widehat u_i=u_i+q(A_i+r).
\]
Then $\widehat c_{i+1}=\widehat c_i+\widehat A_i-
\widehat u_i-\widehat\rho_i$, with
$0\le\widehat\rho_i\le\min(\widehat c_i,\widehat A_i)$ and
$\widehat u_i\ge0$. The identity
\[
 \rho(c+a-u-\rho)=ca-(c-\rho)(a-\rho)-\rho u
 \tag{RD.4}
\]
gives $\widehat\rho_i\widehat c_{i+1}\le
\widehat A_i\widehat c_i$. If the first inequality in (RD.3) fails,
then $\widehat u_i>2\tau\max(\widehat c_i,\widehat A_i)$,
and the bound improves by $1-\tau$: split at
$\widehat\rho_i=\min(\widehat c_i,\widehat A_i)/2$ in (RD.4).
If the second fails, the factor $n_i+r$ instead improves by
$1-\lambda/k+\eta\le e^{-\tau}$.
Telescoping the quadratic inequalities therefore gives
\[
 (\ell_0+r)\prod_i(\rho_i+r)(n_i+r)
 \le\theta A\prod_iA_iC_i,\qquad
 \theta^{r-1}\le
 \exp\!\left(\frac{(k-2)(4s+1)}{k^2}
             -\frac{(k-2)\lambda}{4k}\right)
 \le e^{2-\lambda/5}.
 \tag{RD.5}
\]
The factors lost in regularizing are at most $e^\eta$ initially and
$e^{2\eta}$ per interior degree.

One signed Pieri path to any dominant weight proves
$w_r(\ell)\le\prod_tS_r(|\ell_t|)$: add positive coordinates from
top to bottom, then subtract negative ones from bottom to top, and
apply the total path-dimension identity. The bounds
$y^{r-1}/(r-1)!\le S_r(y)\le(y+r)^{r-1}/(r-1)!$ and (RD.5)
give $w_r(\ell)\le e^{2-\lambda/5}\mathcal D(A)$ for each bad pair.
There are at most $(8a)^{4(N_1+1)^2}$ pairs: every state is determined
by its first $s+1$ and last $s$ coordinates, all bounded by
$A+(2s+1)\lfloor2a/p\rfloor\le2a$.
Even retaining the source's extra envelope
$(8a+k)^{(2N_1+2)^2}$, the total bad weight is at most
\[
 e^{2-\lambda/5+64N_1^2\log a}\mathcal D(A)
 \le e^{-\lambda/8}\mathcal D(A),\qquad C_0\ge880.
 \tag{RD.6}
\]
Thus the kept pairs retain at least $1-e^{-\lambda/8}$ of (RD.2).

\paragraph{Why these pairs give the actual filtered blocks.}
The coarse alternative SW10 says
$A_i(1+\delta)<A_*$ before the split $\sigma$, and
$A_*<A_i(1-\delta)$ after it. From
$c_i\le A+\sum_{l<i}A_l$, the late case gives $c_i\le A_i$ and
\[
                 O_i\le c_i+\lambda A_i/k\le A_i.
\]
In the early case, induction in (RD.3) gives
\[
 c_i\ge A_*(1-e),\quad
 e=4N_1\lambda/k+12N_1/p^2\le9\delta/80,
 \quad e(1+\delta)\le\delta.
\]
Hence $A_i\le c_i$, and
$O_i\le A_i(1+4/p^2)+\lambda c_i/k$,
$S_i\le2\lambda A_*/k$.
Finally, terminal interlacing telescopes to
\[
 D:=\sum\ell-\sum K^{2s+1}
 \le c_s+\sum_iC_i-\sum_i n_i
 \le c_s+(\lambda/k)\sum_iC_i.
\]

For $0\le j_+\le J:=\lfloor\lambda B/k\rfloor$, define
\[
\begin{array}{c|ccc}
 &o_i&w_i&\mathrm{sg}_i\\ \hline
 i<\sigma&0&W_i&S_i\\
 i\ge\sigma&O_i&0&0
\end{array}
 \qquad o_+=D+j_+.
\]
Keep the original digits, choose the late levels $j_i=S_i$, and use
terminal level $j_+$. The preceding inequalities are exactly the
source kept-block inequalities; the terminal one follows also from
$C_i\le2a/P_i$. Moreover
\[
 \mathrm{cap}_l=A_*+\sum_{i<l}(A_i-O_i)=M+c_l.
 \tag{RD.7}
\]
The actual level chain is $K^{n+1}$, followed by $\ell$.
Its degrees are $W_i,C_i-S_i,D$, exactly the level-adjusted block
degrees. Levels are legal because $S_i\le O_i,C_i$ and $j_+\le D+j_+,B$.
At each tested prefix its first coordinate is at most $c_i$, or $c_s$
terminally; (RD.7) and $M\ge h$ give the required
$\lfloor\mathrm{cap}-h\rfloor$ cutoff.
The map is injective: the tail chain and its fixed degree sums recover
every omitted first coordinate; the terminal state recovers $\ell$,
and its level recovers $j_+$. Its weight is
$w_r(\ell)S_r(B-j_+)$. This proves
\[
 \mathrm{chainCount}(h)\ge
 (1-e^{-\lambda/8})\mathcal D(A_*-M)
                  \sum_{j=0}^{J}S_r(B-j).
 \tag{RD.8}
\]
All sums are finite: the kept inequalities bound every block order,
then every level and chain coordinate.

\paragraph{The remaining factors and the half-density window.}
The hockey-stick identity and the product formula for binomials give
\[
\begin{aligned}
 \sum_{j=0}^{J}S_r(B-j)&=S_k(B)-S_k(B-J-1)
                    \ge(1-e^{-\lambda/3})S_k(B),\\
 S_k(A_*-M)&\ge(1-M/A_*)^rS_k(A_*)
                    \ge(1-rM/A_*)S_k(A_*).
\end{aligned}
\]
For the first bound use $J+1>\lambda B/k$ and
$k(B+r)\le3rB$, valid since $r\ge4$ and $B\ge k$.
Applying these to (RD.8), and the actual rank identity RU1--RU2,
proves (RD.1). Its last factor may be negative; that binomial comparison
is used last, with a nonnegative coefficient.

In the source window $0\le h\le b'/(kQY)$,
$MkQ\le2b'$ and $A_*Q\ge a/2\ge60b'$ give
$rM/A_*\le1/30$ and the required $M\le\epsilon A_*$.
Since $\lambda\ge24$, the two exponential factors are at least
$3/4$ and $8/9$. Together with RU8,
\[
       \frac12\,\mathrm{mmB}
       \le\mathrm{rankOnDigits}(h)\le\mathrm{mmB},
       \qquad 0\le h\le\frac{b'}{kQY}.
 \tag{RD.9}
\]
Indeed the proved lower factor is at least $29/45>1/2$.
This is the actual one-segment comparison, including its digit and
rounding conventions. The expanded report proves each intervening
inequality and identifies the exact source declarations.

\subsection{RD11--RD15: harmonic density and the final window mean}
\label{sec:rank-density-mean}

\status{The actual density and column-box numerical records are proved
with the source thresholds. A finite telescoping argument strengthens
the harmonic density and shortens its proof. With HM7 and the completed
local mean and margin bounds, it gives the literal final exceptional
mean $[6(C_A+C_s)+8]NN_1/Y$.}

\paragraph{A finite harmonic calculation.}
In the HM1 convention put $P(t)=\binom{t+r}{r}$, $P(-1)=0$, and
$H(\alpha,\beta)=P(\alpha)P(\beta)-P(\alpha-1)P(\beta-1)$.
Let $R_j=P(\alpha)P(\beta)/[P(\alpha+j)P(\beta+j)]$.
Monotonicity of $t/(t+r)$ gives the exact comparison
\[
 \frac{H(\alpha,\beta)}{P(\alpha+j)P(\beta+j)}
 =R_j\left(1-\frac{\alpha\beta}{(\alpha+r)(\beta+r)}\right)
 \ge R_j-R_{j+1}.
 \tag{RD.10}
\]
Thus its sum for $0\le j\le J$ is at least $1-R_{J+1}$.
Writing $L=J+1$ and discarding the first polynomial ratio,
\[
 R_L\le\prod_{t=1}^{L}\left(1-\frac r{\beta+r+t}\right)
 \le e^{-rL/(\beta+r+L)}\le e^{-\lambda'/3}
 \tag{RD.11}
\]
whenever $r\ge2$, $\beta+L+r\le2b'$ and
$L\ge\lambda'b'/(r+1)$. This improves the source lower density by
removing its extra factor $1-rJ_0/\alpha$; it works also at zero degrees.

\paragraph{The literal enlarged windows.}
Use the final data $X.\mathrm{Hyp}(K.C_0,C_H)$ and assume
$C_H\ge\max(T_{\rm LP},400)$, where NP4 proves the base estimates and
\[
 T_{\rm LP}=\max\{108,8K.C_A+93,1,K.C_0^{-2}\}.
\]
Write $a=\lceil a_{\mathbb R}\rceil$,
$\lambda_0=\mathrm{LaI}(a_{\mathbb R})$,
$\lambda' =\mathrm{LaI}(a)$, $Z=\lambda_0b'/k$,
$J_0=\lambda'b'/k$, $J=\lfloor J_0\rfloor$, $c=c_{\rm cl}$,
and $e=J+c$. The source windows are
\[
\begin{aligned}
 W_a(\tau)&=[a+\tau,\lfloor a_{\mathbb R}+a_{\mathbb R}/(kY)\rfloor+\tau],\\
 W_b(\tau)&=[\lceil b_l\rceil+\tau,\lfloor b_u\rfloor+\tau],
\end{aligned}
\]
intersected with the integers. Their enlarged versions $W_a^+,W_b^+$
raise only the upper endpoints by $e$.
The literal hypotheses and NP4 rounding give
\[
\begin{gathered}
 \lambda_0\le\lambda'\le2\lambda_0,\quad
 \lambda_0\ge8\log(NY),\quad c\le Z,\quad J_0\le2Z,\\
 b'\ge3200k,\quad800Z\le b',\quad Z\ge1,\qquad
 w_a,w_b\ge Y_bNZ/2>0,
\end{gathered}
\]
where $w_a=|W_a(0)|$, $w_b=|W_b(0)|$ and $Y_b\ge400$.
For the window bounds, rounding loses at most one from the real widths;
those widths are at least $Y_bN\lambda_0b'$ and $Y_bNZ$.
For $\tau\le c$ and $\beta\in W_b(\tau)$,
$\beta+J+1+r\le3b'/2+3Z+k<2b'$, so (RD.11) applies uniformly.
Choose $d=1-e^{-\lambda'/3}$. Then
\[
 Ne^{-\lambda'/3}\le\frac13,
 \qquad \frac{Ne}{w_a},\frac{Ne}{w_b}\le\frac3{200},
 \qquad d^N\ge\frac23.
\]
The first inequality follows by cubing:
$N^3e^{-\lambda'}\le N^{-5}\le1/32<1/27$.
Bernoulli also gives
$ (1+u)^N\le(1-Nu)^{-1}$ for $Nu<1$.
Consequently
\[
 |W_a^+|^N|W_b^+|^N
 \le(200/197)^2w_a^Nw_b^N
 \le2d^Nw_a^Nw_b^N.
 \tag{RD.12}
\]
Together with (RD.10)--(RD.11), this proves every field of the actual
$\mathrm{DensityNumericalFacts}$. HM7 therefore gives the two translated
averages needed below.

\paragraph{Boundary loss and the exceptional mean.}
Already $C_H\ge T_{\rm LP}$ implies
$w_a,w_b\ge2YNc$ and $12Nkc/b'\le1/Y$.
Indeed the size clause gives $C_Hk^3NYc\le b'$,
while the window widths satisfy
$w_a\ge216b'-1$ and $w_b\ge216b'/k-1$; also $b'\ge k$.
Bernoulli now bounds the fraction below a shifted lower endpoint by
\[
 1-\left(1-\frac c{w_a}\right)^N
       \left(1-\frac c{w_b}\right)^N\le\frac1Y.
 \tag{RD.13}
\]
These are exactly the two nontrivial fields of
$\mathrm{ColumnSpanNumericalFacts}$.

SW8--SW9, RU8 and NP3--NP6 give each actual segment mean at most
$C_A(s_j+1)N_1/Y$; NP9 gives $\kappa_P\le C_sNN_1/Y$.
The sufficient fixed coefficients are
\[
 C_A\ge\max\{56,2(31e^3+5)e^3,1\},\qquad
 C_s\ge C_\kappa(C_m+4C_{E1}+C_{E2}+15).
\]
Each assignment on the $s_j+1$ distinct segment coordinates has
$|W_a^+|^{N-s_j-1}$ extensions. Finite marginalization, followed by
$\sum_j(s_j+1)\le N$, therefore proves
$\operatorname{mean}k_x\le(C_A+C_s)NN_1/Y$.
No independence between different segments is needed.

The source $k_{\rm text}$ is the inside term $k_F(A-c,B-c)$
(replaced by $1$ outside), plus $2k_F(A+c,B+c)+12Nkc/b'$.
Subtraction is injective on the inside box; addition is a bijection
onto the translated box. Since $k_F\ge0$, HM7 and (RD.13) give
\[
 \operatorname{mean}k_{\rm text}
 \le6\operatorname{mean}k_x+\frac2Y
 \le[6(C_A+C_s)+8]\frac{NN_1}{Y}.
 \tag{RD.14}
\]
The coefficient is exactly $K.C_{\rm mean}$. The expanded report records
the natural-to-signed window identification and all constant thresholds.
